\section{The specified whole-section regularized operator}
\textit{Supports: G331, n1, n2.}

Let $\phi:X_R\to[-\infty,+\infty)$ be USC and bounded above, with partial bottom allowed. Define, on the ENTIRE compact section,
\[
p_R\phi=\sup\{u\in D_R:u\le\phi\},\qquad E_R\phi=(p_R\phi)^*.
\]
Bottom is a competitor. Nonpositivity and the USC obstacle give
\[
p_R\phi\le E_R\phi\le\min(\phi,0).                    \tag{6}
\]
There is no positivity assumption in the definition. Exactly
\[
E_R\phi\in D_R\ \Longleftrightarrow\ \text{the competitors have an actual greatest }D_R\text{ member}. \tag{7}
\]
If positive, (6) makes $E$ a competitor, hence $E=p$. Conversely a greatest member equals $p$, is USC, and therefore equals its regularization. This includes the all-bottom case.

For finite continuous $f$, the raw suprema using $C$ and $D$ are equal. If $u\in D$ is below $f$, an actual representation $g_\alpha\downarrow u$ has, for each $\epsilon>0$, one term $g_\alpha<f+\epsilon$ everywhere: continuous strict-test open sets cover the compact section, and directedness selects a common index. Then $g_\alpha-\epsilon\in C$ is below $f$ and above $u-\epsilon$. Let $\epsilon$ decrease. This does NOT claim that every USC obstacle has this reduction. The common raw supremum $p$ is LSC and finite, since a sufficiently negative constant is below $f$. Thus
\[
E_Rf\in D_R\ \Longleftrightarrow\ p_Rf\text{ is continuous},
\quad\text{and then } E_Rf=p_Rf\in C_R.                \tag{8}
\]
For the forward implication use (7), which makes the LSC $p$ USC. For the reverse, choose a $C$ competitor within $\epsilon$ of $p$ at each point; continuity of $p$ and those competitors gives open neighborhoods. A finite subcover and finite maxima yield a $C$ competitor uniformly within $\epsilon$ of $p$. Uniform closure proves $p\in C$. On a reduced curve the granted geometric class is the full nonpositive branch orthant, so $E\phi=\min(\phi,0)$ on the finite section. Empty sections have their unique empty function. Equations (7)--(8) give a selected positive domain, not universal existence of positive envelopes. At the qualified smooth BFJ scope, G331 identifies $D$ with the full canonical BFJ all-semivaluation class with bottom adjoined. Consequently this is exactly the USC supremum of that SAME class; any smooth envelope with this signature and an actual greatest competitor agrees by (7). A published fixed-skeleton nef-divisor envelope has a narrower obstacle signature; no assertion of general smooth USC-obstacle existence is inferred from it.

For either map class in n2, and every such extended USC obstacle,
\[
p_B(\phi\circ\rho)=p_R\phi\circ\rho.                \tag{9}
\]
One inequality is pullback. For the other, take a source $D$ competitor; its full fibre maximum is a downstairs $D$ function below $\phi$, dominating that competitor on each lift. This proves the raw equality for all supports and partial-bottom patterns. The map is open by n1--n2, so (2) gives the actual SPECIFIED regularized identity
\[
E_B(\phi\circ\rho)=E_R\phi\circ\rho.                \tag{10}
\]
This is an equality of functions even outside the positive domain. Full positive descent proves exact equivalence of their positive domains; for finite continuous obstacles (8) also identifies both positive outputs with $C$. We have used positive maxima AND openness, not merely ideal extrema or a raw supremum. No normal-target assertion with vertical components is made.

For an arbitrary reduced finite INJECTION define the constrained upstairs envelope by taking only FULL-fibre-constant $D_B$ competitors below $\phi\circ\rho$, then ordinary upstairs USC regularization; denote it $E_{B,\mathrm{fib}}$. G331's completed positive descent identifies its raw supremum with $p_R\phi\circ\rho$. Applying compact quotient formula (1) proves the exact general rule
\[
E_R\phi(x)=\max_{\rho y=x}E_{B,\mathrm{fib}}(\phi\circ\rho)(y). \tag{11}
\]
In particular downstairs positivity holds exactly when $E_{B,\mathrm{fib}}$ is BOTH in $D_B$ AND constant on EVERY fibre. Indeed positive downstairs output equals the raw supremum by (7), whose pullback is USC and makes the constrained regularization equal to it. Conversely a positive fibre-constant constrained regularization descends positively by G331; it is below $\phi$ and above $p_R\phi$, so equality follows. The fibre condition cannot be dropped after regularization.

The unconstrained source envelope dominates the constrained one. Thus the universal comparison is
\[
\rho^*E_R\phi\le M_{\mathrm{fib}} E_B(\phi\circ\rho),  \tag{12}
\]
where $M_{\mathrm{fib}}H(y)=\max_{\rho z=\rho y}H(z)$. Without openness this does not imply $\rho^*E_R\phi\le E_B$. If $E_B$ is positive, equality $E_B=\rho^*E_R$ holds exactly when $E_B$ is FULL-fibre constant: descent of its greatest competitor proves sufficiency, and necessity is immediate. If both envelopes are positive, pulled downstairs output is a source competitor, giving that pointwise inequality; equality still requires full constancy. For normalization, ``full'' means the entire conductor quotient, including nonminimal supports and all lower-dimensional points. These are operator statements, rather than definitions of another order completion.
